\chapter{Có Ai Thấy Em Đang Kiệt Sức Không}
\markboth{Có Ai Thấy Em Đang Kiệt Sức Không}{Does Anyone Notice I Am Running Out}

\section{Em Vẫn Đi Học Đều Nên Mọi Người Tưởng Em Ổn}
\begin{truyen}{Em Vẫn Đi Học Đều Nên Mọi Người Tưởng Em Ổn}{I Still Show Up, So Everyone Thinks I Am Fine}
\chuhoa{C}{ó những học sinh chưa nghỉ buổi nào, vẫn có mặt đầy đủ, vẫn nộp bài tương đối đúng.}
\textit{(Some students have perfect attendance, still show up daily, and submit work more or less on time.)}

Chính vì còn làm được như vậy nên người lớn nghĩ em ổn.
\textit{(Precisely because they still function, adults assume they are fine.)}

Nhưng kiệt sức không phải lúc nào cũng đến bằng sự sụp hẳn.
\textit{(Burnout does not always arrive as collapse.)}

Có em nhìn xuống bàn rồi nói: “Em vẫn đi học được, nhưng em mệt kiểu hết sạch bên trong.”
\textit{(One student looked down and said, “I can still come to school, but inside I feel completely used up.”)}

Nó có thể đến bằng việc một đứa trẻ vẫn gồng được thêm vài tuần nữa.
\textit{(It may arrive through a child managing to hold on for a few more weeks.)}

\end{truyen}

\begin{baihoc}
Một học sinh còn đang cố chưa chắc là một học sinh còn đang ổn.
\textit{(A student who is still trying is not necessarily a student who is still okay.)}
\end{baihoc}

\begin{ghinhoanh}
\textbf{Short English to remember:}

Still functioning does not always mean still okay.

\end{ghinhoanh}

\begin{tuhocnhanh}
\textbf{Gợi ý để dạy và tự nghĩ / Teaching and reflection prompts}

1. Vì sao sự đều đặn dễ làm người lớn chủ quan? \textit{Why does consistency make adults assume too much?}

2. Kiệt sức kín đáo trông như thế nào? \textit{What does quiet exhaustion look like?}

3. Mình có đang nhầm gồng với khỏe không? \textit{Am I mistaking endurance for health?}
\end{tuhocnhanh}
\ngancach

\section{Em Nói Không Sao Vì Không Muốn Phiền}
\begin{truyen}{Em Nói Không Sao Vì Không Muốn Phiền}{I Say I Am Fine Because I Do Not Want to Trouble Anyone}
\chuhoa{C}{ó em được hỏi thì luôn đáp rất gọn: em không sao ạ.}
\textit{(Some students answer every concern with a short line: I am fine.)}

Câu đó nhiều khi không phải sự thật trọn vẹn.
\textit{(That sentence is often not the whole truth.)}

Nó là cách em tiết kiệm người lớn, tiết kiệm thời gian, và tiết kiệm cả sự thất vọng nếu nói ra mà không ai hiểu.
\textit{(It is a way of sparing adults, saving time, and avoiding the disappointment of not being understood.)}

Có những học sinh rất mệt vì quá quen với việc tự gói phần mình lại cho gọn.
\textit{(Some students are exhausted because they are too used to folding their pain into a small package.)}

\end{truyen}

\begin{baihoc}
Câu em không sao đôi khi là lời nói tiện nhất, không phải lời nói thật nhất.
\textit{(I am fine is sometimes the most convenient answer, not the truest one.)}
\end{baihoc}

\begin{ghinhoanh}
\textbf{Short English to remember:}

I am fine can be the easiest answer, not the truest one.

\end{ghinhoanh}

\begin{tuhocnhanh}
\textbf{Gợi ý để dạy và tự nghĩ / Teaching and reflection prompts}

1. Vì sao học sinh chọn câu không sao? \textit{Why do students choose the line I am fine?}

2. Câu này có thể che điều gì? \textit{What can this answer be hiding?}

3. Người dạy nên hỏi tiếp theo cách nào? \textit{How should a teacher continue the conversation?}
\end{tuhocnhanh}
\ngancach

\section{Em Không Lười, Em Chỉ Cạn Pin}
\begin{truyen}{Em Không Lười, Em Chỉ Cạn Pin}{I Am Not Lazy, I Am Depleted}
\chuhoa{C}{ó những em đầu tuần còn cố được, càng về cuối tuần càng đuối hẳn.}
\textit{(Some students manage early in the week and then fade badly by the end of it.)}

Người lớn nhìn vào thường chỉ thấy sa sút.
\textit{(Adults often only see decline.)}

Nhưng có những thân thể và tinh thần đã đi quá tải lâu rồi.
\textit{(But some bodies and minds have been overloaded for too long.)}

Nếu cứ gọi đó là lười, mình sẽ bỏ lỡ đúng chỗ cần cứu nhất.
\textit{(If we keep calling it laziness, we miss the exact place that needs help most.)}

\end{truyen}

\begin{baihoc}
Không phải mọi sự chậm lại đều là thiếu ý chí. Có cái là dấu hiệu của cạn kiệt.
\textit{(Not every slowing down is a lack of will. Some of it is a sign of depletion.)}
\end{baihoc}

\begin{ghinhoanh}
\textbf{Short English to remember:}

Not every slowdown is laziness. Some of it is depletion.

\end{ghinhoanh}

\begin{tuhocnhanh}
\textbf{Gợi ý để dạy và tự nghĩ / Teaching and reflection prompts}

1. Vì sao cạn kiệt dễ bị gọi là lười? \textit{Why is depletion often mislabeled as laziness?}

2. Dấu hiệu nào cho thấy đây là quá tải? \textit{What signs suggest overload instead?}

3. Mình đang dùng nhãn nào quá nhanh? \textit{What label am I applying too quickly?}
\end{tuhocnhanh}
\ngancach

\section{Em Cần Nghỉ Một Chút Có Được Không}
\begin{truyen}{Em Cần Nghỉ Một Chút Có Được Không}{Am I Allowed a Moment of Rest}
\chuhoa{C}{ó những học sinh không cần bỏ cuộc. Các em chỉ cần một quãng thở ra, một nhịp chậm lại, một ngày bớt bị đòi.}
\textit{(Some students do not need to quit. They only need room to breathe, a slower pace, one day of being demanded less.)}

Nhưng nhiều môi trường học đường rất khó cho phép chuyện đó.
\textit{(But many school environments find that hard to allow.)}

Người lớn sợ nghỉ là trượt.
\textit{(Adults fear that rest means falling behind.)}

Trong khi có những trường hợp không được nghỉ mới là thứ làm em trượt thật.
\textit{(Yet in some cases, the lack of rest is exactly what causes the real fall.)}

\end{truyen}

\begin{baihoc}
Nghỉ đúng lúc không phải phản bội sự cố gắng. Nhiều khi nó cứu sự cố gắng khỏi gãy hẳn.
\textit{(Resting at the right time does not betray effort. It often saves effort from breaking completely.)}
\end{baihoc}

\begin{ghinhoanh}
\textbf{Short English to remember:}

Timely rest can save effort from breaking.

\end{ghinhoanh}

\begin{tuhocnhanh}
\textbf{Gợi ý để dạy và tự nghĩ / Teaching and reflection prompts}

1. Vì sao người lớn hay sợ việc nghỉ? \textit{Why are adults afraid of rest?}

2. Nghỉ đúng lúc giúp được gì? \textit{What can timely rest protect?}

3. Mình có đang để học sinh kiệt mới cho dừng không? \textit{Am I waiting until collapse before allowing pause?}
\end{tuhocnhanh}
\ngancach

\section{Em Ngồi Đó Nhưng Gần Như Hết Pin Rồi}
\begin{truyen}{Em Ngồi Đó Nhưng Gần Như Hết Pin Rồi}{I Am Sitting Here, But My Battery Is Nearly Gone}
\chuhoa{C}{ó em vẫn mở vở, vẫn cầm bút, vẫn nhìn bảng. Nhưng nhìn kỹ là thấy phần tỉnh táo của em gần như chạm đáy.}
\textit{(Some students still open the notebook, hold the pen, and face the board. But a closer look shows that their alertness is nearly gone.)}

Các em chưa gục hẳn nên dễ bị bỏ qua.
\textit{(They have not collapsed, so they are easy to overlook.)}

Nhiều sự kiệt sức nguy hiểm nằm đúng ở vùng chưa sập nhưng đã cạn.
\textit{(Many dangerous forms of exhaustion live exactly in that zone: not fallen, but emptied.)}

Người dạy tinh là người nhận ra phần cạn đó trước khi học sinh rơi hẳn.
\textit{(A perceptive teacher notices that emptiness before the student fully falls.)}

\end{truyen}

\begin{baihoc}
Học sinh chưa gục không có nghĩa là học sinh còn đủ sức.
\textit{(A student who has not collapsed is not necessarily a student who still has strength.)}
\end{baihoc}

\begin{ghinhoanh}
\textbf{Short English to remember:}

Not collapsed does not mean not exhausted.

\end{ghinhoanh}

\begin{tuhocnhanh}
\textbf{Gợi ý để dạy và tự nghĩ / Teaching and reflection prompts}

1. Kiệt sức trước khi sụp hẳn trông ra sao? \textit{What does exhaustion look like before collapse?}

2. Vì sao vùng “chưa gục nhưng đã cạn” dễ bị bỏ qua? \textit{Why is the zone of “not collapsed but emptied” easy to miss?}

3. Mình có nhìn đủ kỹ phần năng lượng của học sinh không? \textit{Am I observing students' energy carefully enough?}
\end{tuhocnhanh}
\ngancach

\section{Em Xin Ra Ngoài Không Phải Để Trốn Học}
\begin{truyen}{Em Xin Ra Ngoài Không Phải Để Trốn Học}{I Ask to Step Out, Not to Escape the Lesson}
\chuhoa{C}{ó em trong một tiết xin ra ngoài tới hai lần. Nhìn nhanh rất dễ nghĩ là né học.}
\textit{(Some students ask to step out twice in one period. A quick glance makes it look like avoidance.)}

Nhưng ra hành lang, em chỉ đứng tựa tường, hít rất sâu.
\textit{(But in the hallway the student only leans on the wall and breathes deeply.)}

Tôi hỏi: “Em ổn không?”
\textit{(I asked, “Are you okay?”)}

Em nói: “Em chỉ muốn đầu em đừng ồn nữa.”
\textit{(The student said, “I just want my head to stop being noisy.”)}

\end{truyen}

\begin{baihoc}
Có những lúc học sinh ra khỏi lớp không phải để bỏ học, mà để giữ mình khỏi vỡ thêm.
\textit{(Sometimes students leave the room not to escape learning, but to keep themselves from breaking further.)}
\end{baihoc}

\begin{ghinhoanh}
\textbf{Short English to remember:}

Sometimes stepping out is a way of holding oneself together.

\end{ghinhoanh}

\begin{tuhocnhanh}
\textbf{Gợi ý để dạy và tự nghĩ / Teaching and reflection prompts}

1. Vì sao hành vi xin ra ngoài dễ bị hiểu sai? \textit{Why is this behavior easily misunderstood?}

2. Học sinh đang cố giữ điều gì? \textit{What is the student trying to hold together?}

3. Mình có thể quan sát thêm gì trước khi kết luận? \textit{What can I observe before drawing conclusions?}
\end{tuhocnhanh}
\ngancach

\section{Em Không Muốn Cố Thêm Nữa, Em Muốn Được Thấy Em Đang Hết Sức}
\begin{truyen}{Em Không Muốn Cố Thêm Nữa, Em Muốn Được Thấy Em Đang Hết Sức}{I Do Not Need More Pressure, I Need Someone to See I Am Done}
\chuhoa{C}{ó em bị động viên suốt: cố lên, ráng lên, thêm chút nữa.}
\textit{(Some students are constantly encouraged: try harder, keep going, just a little more.)}

Những câu đó không sai, nhưng có lúc đi trật thời điểm.
\textit{(Those lines are not wrong, but sometimes they arrive at the wrong moment.)}

Khi học sinh đã chạm đáy, điều em cần không phải một cú thúc mới.
\textit{(When a student has hit the bottom, what they need is not another push.)}

Em cần có người nhìn thẳng và nói: “Ừ, thầy thấy em đang hết sức rồi.”
\textit{(They need someone to look directly at them and say, “Yes, I can see you are at your limit.”)}

\end{truyen}

\begin{baihoc}
Được công nhận là đã cố hết sức nhiều khi chữa được nhiều hơn bị thúc thêm.
\textit{(Being recognized as already trying one's best can heal more than more pressure.)}
\end{baihoc}

\begin{ghinhoanh}
\textbf{Short English to remember:}

Sometimes students need recognition of effort, not more pressure.

\end{ghinhoanh}

\begin{tuhocnhanh}
\textbf{Gợi ý để dạy và tự nghĩ / Teaching and reflection prompts}

1. Vì sao lời động viên đúng cũng có thể sai lúc? \textit{Why can even a good encouragement be wrongly timed?}

2. Học sinh cần nghe gì ở đáy sức? \textit{What do students need to hear at the limit of their strength?}

3. Mình có đang thúc thêm khi học sinh cần được nhìn ra hơn không? \textit{Am I pushing when students need to be seen instead?}
\end{tuhocnhanh}
\ngancach

\section{Tối Nào Em Cũng Muộn Nên Sáng Nào Em Cũng Đuối}
\begin{truyen}{Tối Nào Em Cũng Muộn Nên Sáng Nào Em Cũng Đuối}{Every Late Night Becomes Another Heavy Morning}
\chuhoa{C}{ó em sáng nào cũng ngáp, mắt sưng, đầu chậm.}
\textit{(Some students arrive every morning yawning, puffy-eyed, and mentally slow.)}

Hỏi ra mới biết tối nào cũng thức rất khuya vì học, vì phụ việc, hoặc vì đầu óc không yên nổi.
\textit{(When asked, they reveal that every night they stay up late to study, help at home, or simply because their minds cannot settle.)}

Ban ngày bị trách thiếu tập trung.
\textit{(In the daytime they are blamed for poor concentration.)}

Nhưng có những buổi sáng mệt là hóa đơn của những đêm kéo quá dài.
\textit{(Yet some tired mornings are the bill paid for nights that were stretched too long.)}

\end{truyen}

\begin{baihoc}
Muốn hiểu sự đuối của học sinh, đôi khi phải nhìn vào đêm của các em chứ không chỉ nhìn buổi sáng trên lớp.
\textit{(To understand student fatigue, adults sometimes must look at their nights, not only their mornings in class.)}
\end{baihoc}

\begin{ghinhoanh}
\textbf{Short English to remember:}

Some tired mornings are the bill for nights that were too long.

\end{ghinhoanh}

\begin{tuhocnhanh}
\textbf{Gợi ý để dạy và tự nghĩ / Teaching and reflection prompts}

1. Điều gì đang kéo đêm của học sinh quá dài? \textit{What is stretching students' nights too long?}

2. Vì sao chỉ nhìn biểu hiện trên lớp là chưa đủ? \textit{Why is looking only at classroom behavior not enough?}

3. Mình có đang hiểu sự mệt của học sinh theo chiều 24 giờ không? \textit{Am I understanding student fatigue across the whole day?}
\end{tuhocnhanh}
\ngancach

\section{Một Câu “Hôm Nay Em Mệt Lắm Hả” Có Thể Làm Em Muốn Khóc}
\begin{truyen}{Một Câu “Hôm Nay Em Mệt Lắm Hả” Có Thể Làm Em Muốn Khóc}{One Gentle Question Can Bring Tears}
\chuhoa{C}{ó em cả buổi không nói gì. Đến cuối giờ tôi chỉ hỏi: “Hôm nay em mệt lắm hả?”}
\textit{(A student said nothing all lesson. At the end I only asked, “Are you very tired today?”)}

Em gật đầu rồi khóc.
\textit{(The student nodded and cried.)}

Không phải vì câu hỏi quá hay.
\textit{(Not because the question was brilliant.)}

Chỉ vì có những mệt mỏi đã nằm đó lâu và cuối cùng gặp được một chỗ để rơi xuống.
\textit{(Simply because some exhaustion had been sitting there for a long time and finally found a place to fall.)}

\end{truyen}

\begin{baihoc}
Nhiều lúc học sinh không cần lời giải. Các em chỉ cần một câu hỏi đủ mềm để được vỡ ra.
\textit{(Many times students do not need a solution. They only need a question soft enough to let them break open.)}
\end{baihoc}

\begin{ghinhoanh}
\textbf{Short English to remember:}

Sometimes one gentle question gives exhaustion a place to fall.

\end{ghinhoanh}

\begin{tuhocnhanh}
\textbf{Gợi ý để dạy và tự nghĩ / Teaching and reflection prompts}

1. Vì sao một câu hỏi nhẹ lại làm học sinh bật khóc? \textit{Why can a gentle question bring a student to tears?}

2. Điều gì đã tích lại trước khoảnh khắc đó? \textit{What had been building up before that moment?}

3. Mình có còn giữ được kiểu hỏi mềm như vậy không? \textit{Can I still ask that softly?}
\end{tuhocnhanh}
\ngancach

\section{Em Không Muốn Bỏ Học, Em Chỉ Muốn Bớt Kiệt}
\begin{truyen}{Em Không Muốn Bỏ Học, Em Chỉ Muốn Bớt Kiệt}{I Do Not Want to Quit School, I Just Want Not to Be So Exhausted}
\chuhoa{C}{ó những học sinh nói câu rất thật: “Em không muốn nghỉ học đâu, em chỉ mệt quá.”}
\textit{(Some students say something very honest: “I don't want to quit school. I'm just too tired.”)}

Người lớn nghe chữ mệt hay tưởng các em yếu.
\textit{(Adults hear the word tired and assume weakness.)}

Nhưng nhiều khi đó là một lời kêu cứu rất tỉnh.
\textit{(But often it is a very lucid cry for help.)}

Học sinh còn muốn đi tiếp, chỉ là không muốn bị ép đi tiếp bằng phần pin đã cạn rồi.
\textit{(The student still wants to continue, just not on a battery that is already empty.)}

\end{truyen}

\begin{baihoc}
Khi học sinh nói em mệt, người lớn nên nghe đó như một dữ liệu quan trọng chứ không phải một lời yếu đuối.
\textit{(When students say they are tired, adults should hear it as important information, not as weakness.)}
\end{baihoc}

\begin{ghinhoanh}
\textbf{Short English to remember:}

Student fatigue is important information, not a weak excuse.

\end{ghinhoanh}

\begin{tuhocnhanh}
\textbf{Gợi ý để dạy và tự nghĩ / Teaching and reflection prompts}

1. Vì sao câu em mệt cần được nghe nghiêm túc? \textit{Why should “I am tired” be taken seriously?}

2. Học sinh đang xin điều gì khi nói như vậy? \textit{What are students asking for when they say this?}

3. Mình có đang nghe chữ mệt như dữ liệu giáo dục chưa? \textit{Am I hearing tiredness as educational data yet?}
\end{tuhocnhanh}
